\section*{Capítulo \ref{ch:g10:electron-shells} --- Camadas eletrônicas e tabela periódica}

\begin{solution}{exo:g10:electron-shells:1}
C: $1\mathrm{s}^2\,2\mathrm{s}^2\,2\mathrm{p}^2$. N:
$1\mathrm{s}^2\,2\mathrm{s}^2\,2\mathrm{p}^3$. Mg:
$1\mathrm{s}^2\,2\mathrm{s}^2\,2\mathrm{p}^6\,3\mathrm{s}^2$.
\end{solution}

\begin{solution}{exo:g10:electron-shells:2}
Carbono 4, nitrogênio 5, magnésio 2.
\end{solution}

\begin{solution}{exo:g10:electron-shells:3}
\omterm{def:g10:electron-shells:configuration}{Camada} externa $n = 3$: período 3. $2 + 3 = 5$ \omterm{def:g10:electron-shells:valence}{elétrons de valência}:
\omterm{def:g9:periodic-table-first-look:period}{grupo} 15 (fósforo).
\end{solution}

\begin{solution}{exo:g10:electron-shells:4}
\omterm{def:g10:electron-shells:family}{Metais alcalinos} (lítio, sódio, potássio); \omterm{def:g10:electron-shells:family}{halogênios} (flúor, cloro,
bromo, iodo); \omterm{def:g10:electron-shells:family}{gases nobres} (hélio, neônio, argônio).
\end{solution}

\begin{solution}{exo:g10:electron-shells:5}
Os \omterm{def:g7:atoms-and-molecules:atom}{átomos} tendem a ganhar ou perder \omterm{def:g9:inside-the-atom:nucleus}{elétrons} de modo a ter oito \omterm{def:g9:inside-the-atom:nucleus}{elétrons}
na \omterm{def:g10:electron-shells:configuration}{camada} externa, como o \omterm{def:g10:electron-shells:family}{gás nobre} mais próximo.
\end{solution}

\begin{solution}{exo:g10:electron-shells:6}
K, $2.8.8.1$, perde um: \ce{K+},
$1\mathrm{s}^2\,2\mathrm{s}^2\,2\mathrm{p}^6\,3\mathrm{s}^2\,3\mathrm{p}^6$
(argônio). F, $2.7$, ganha um: \ce{F-},
$1\mathrm{s}^2\,2\mathrm{s}^2\,2\mathrm{p}^6$ (neônio).
\end{solution}

\begin{solution}{exo:g10:electron-shells:7}
$1\mathrm{s}^2\,2\mathrm{s}^2\,2\mathrm{p}^6\,3\mathrm{s}^2$: $Z = 12$,
magnésio.
\end{solution}

\begin{solution}{exo:g10:electron-shells:8}
\ce{Mg^{2+}} e \ce{Cl-}: \ce{MgCl2}. \ce{Al^{3+}} e \ce{O^{2-}}:
\ce{Al2O3}.
\end{solution}

\begin{solution}{exo:g10:electron-shells:9}
O enxofre (6 \omterm{def:g10:electron-shells:valence}{elétrons de valência}, $2.8.6$), no mesmo grupo, o 16.
\end{solution}

\begin{solution}{exo:g10:electron-shells:10}
A \omterm{def:g10:electron-shells:configuration}{camada} externa dele já está completa (oito \omterm{def:g9:inside-the-atom:nucleus}{elétrons}): ele não tem
nada a ganhar perdendo ou ganhando \omterm{def:g9:inside-the-atom:nucleus}{elétrons}.
\end{solution}

\begin{solution}{exo:g10:electron-shells:11}
O argônio tem 18 \omterm{def:g9:inside-the-atom:nucleus}{elétrons}; \ce{X^{2+}} perdeu 2, então X tem 20: o cálcio.
\end{solution}

\begin{solution}{exo:g10:electron-shells:12}
O hélio, $1\mathrm{s}^2$, tem a \omterm{def:g10:electron-shells:configuration}{camada} externa completa (a \omterm{def:g10:electron-shells:configuration}{camada} 1 só
comporta dois \omterm{def:g9:inside-the-atom:nucleus}{elétrons}), como os outros \omterm{def:g10:electron-shells:family}{gases nobres}; ele não se
comporta como os \omterm{def:g9:periodic-table-first-look:metal}{metais} do \omterm{def:g9:periodic-table-first-look:period}{grupo} 2.
\end{solution}

\begin{solution}{exo:g10:electron-shells:13}
O \omterm{def:g10:electron-shells:valence}{elétron de valência} dele está na \omterm{def:g10:electron-shells:configuration}{camada} 4, mais longe do \omterm{def:g9:inside-the-atom:nucleus}{núcleo} que o
do sódio, na \omterm{def:g10:electron-shells:configuration}{camada} 3: menos preso, ele é perdido com mais facilidade.
\end{solution}

\begin{solution}{exo:g10:electron-shells:14}
Os quatro têm 18 \omterm{def:g9:inside-the-atom:nucleus}{elétrons} e a configuração do argônio:
\[ 1\mathrm{s}^2\,2\mathrm{s}^2\,2\mathrm{p}^6\,3\mathrm{s}^2\,3\mathrm{p}^6. \]
\end{solution}

\begin{solution}{exo:g10:electron-shells:15}
Ele ganha 3 \omterm{def:g9:inside-the-atom:nucleus}{elétrons} para alcançar um octeto: tem 5 \omterm{def:g10:electron-shells:valence}{elétrons de
valência}, \omterm{def:g9:periodic-table-first-look:period}{grupo} 15. No período 3, é o fósforo:
\[ 1\mathrm{s}^2\,2\mathrm{s}^2\,2\mathrm{p}^6\,3\mathrm{s}^2\,3\mathrm{p}^3. \]
\end{solution}

\begin{solution}{pb:g10:electron-shells:1}
\textbf{1.} Alumínio: 13 \omterm{def:g9:inside-the-atom:proton}{prótons}, 13 \omterm{def:g9:inside-the-atom:nucleus}{elétrons}. Oxigênio: 8 \omterm{def:g9:inside-the-atom:proton}{prótons}, 8
\omterm{def:g9:inside-the-atom:nucleus}{elétrons}.

\textbf{2.} Al: $1\mathrm{s}^2\,2\mathrm{s}^2\,2\mathrm{p}^6\,3\mathrm{s}^2\,3\mathrm{p}^1$.
O: $1\mathrm{s}^2\,2\mathrm{s}^2\,2\mathrm{p}^4$.

\textbf{3.} Al: 3 \omterm{def:g10:electron-shells:valence}{elétrons de valência}, período 3, \omterm{def:g9:periodic-table-first-look:period}{grupo} 13. O: 6,
período 2, \omterm{def:g9:periodic-table-first-look:period}{grupo} 16.

\textbf{4.} O alumínio é um \omterm{def:g9:periodic-table-first-look:metal}{metal}, o oxigênio um \omterm{def:g9:periodic-table-first-look:metal}{não metal}.

\textbf{5.} \ce{Al^{3+}}: $1\mathrm{s}^2\,2\mathrm{s}^2\,2\mathrm{p}^6$.

\textbf{6.} \ce{O^{2-}}: $1\mathrm{s}^2\,2\mathrm{s}^2\,2\mathrm{p}^6$.

\textbf{7.} O neônio.

\textbf{8.} O alumínio perde 3, o oxigênio ganha 2.

\textbf{9.} $2 \times (+3) + 3 \times (-2) = 0$: \ce{Al2O3}.

\textbf{10.} Perdidos: $2 \times 3 = 6$. Ganhos: $3 \times 2 = 6$.

\textbf{11.} Os \omterm{def:g9:inside-the-atom:nucleus}{elétrons} não são criados nem destruídos: os que o
alumínio perde são exatamente os que o oxigênio ganha, e o composto é
neutro.

\textbf{12.} Ele tem a mesma carga, $+3$: as cargas do cristal continuam
equilibradas.

\textbf{13.} $2 \times 27 + 3 \times 16 = 102$ \omterm{def:g9:inside-the-atom:proton}{núcleons}.

\textbf{14.} $102 \times 1.67 \times 10^{-27} = \qty{1.70e-25}{kg}$.

\textbf{15.} $10^{-3} / 1.70 \times 10^{-25} \approx 5.87 \times 10^{21}$
unidades de fórmula.

\textbf{16.} $2 \times 5.87 \times 10^{21} = 1.17 \times 10^{22}$
\omterm{def:g9:ions:ion}{íons} \ce{Al^{3+}}; $3 \times 5.87 \times 10^{21} = 1.76 \times 10^{22}$
\omterm{def:g9:ions:ion}{íons} \ce{O^{2-}}.

\textbf{17.} Um \omterm{def:g9:inside-the-atom:nucleus}{elétron} é cerca de 1836 vezes mais leve que um \omterm{def:g9:inside-the-atom:proton}{núcleon};
os 50 \omterm{def:g9:inside-the-atom:nucleus}{elétrons} de uma unidade de fórmula pesam menos de três centésimos
de um \omterm{def:g9:inside-the-atom:proton}{núcleon}.

\textbf{18.} $6 \times 5.87 \times 10^{21} \approx 3.5 \times 10^{22}$
\omterm{def:g9:inside-the-atom:nucleus}{elétrons} transferidos para \qty{1}{g} de coríndon.
\end{solution}
